\section*{Capítulo \ref{ch:g12:human-evolution} --- Um olhar sobre a evolução humana}

\begin{solution}{exo:g12:human-evolution:1}
O chimpanzé (com o bonobo); as linhagens se separaram seis a sete
milhões de anos atrás.
\end{solution}

\begin{solution}{exo:g12:human-evolution:2}
Abertura medular embaixo do crânio; uma coluna de curva dupla; uma
bacia curta em forma de tigela; fêmures inclinados para dentro até o
joelho; um pé arqueado com o dedão alinhado.
\end{solution}

\begin{solution}{exo:g12:human-evolution:3}
Cerca de \qty{420}{cm^3}; cerca de \qty{850}{cm^3}; cerca de \qty{1350}{cm^3}.
\end{solution}

\begin{solution}{exo:g12:human-evolution:4}
Porque várias \omterm{def:g10:biodiversity-scales:species}{espécies} da linhagem viveram ao mesmo tempo, e a maioria
delas terminou em extinção em vez de virar a seguinte; as \omterm{def:g10:biodiversity-scales:species}{espécies} são
ramos, não degraus.
\end{solution}

\begin{solution}{exo:g12:human-evolution:5}
Na África, há cerca de 300\,000 anos. A \omterm{def:g10:biodiversity-scales:biodiversity}{diversidade genética} é maior na
África e cai com a distância a partir dela: todos os seres humanos
vivos descendem daquela \omterm{def:g12:selection-drift-speciation:population}{população} africana, através de grupos
fundadores pequenos e sucessivos.
\end{solution}

\begin{solution}{exo:g12:human-evolution:6}
Dois \omterm{def:g10:universal-dna:information}{cromossomos} ancestrais se fundiram na \omterm{def:g12:human-evolution:lineage}{linhagem humana} num só, e
assim 24 pares viraram 23. O \omterm{def:g11:cell-cycle-mitosis:chromosome}{cromossomo} fundido carrega, no meio, os
restos das duas pontas originais e um segundo \omterm{def:g11:cell-cycle-mitosis:chromosome}{centrômero}, silencioso
--- um registro da fusão; o conteúdo de \omterm{def:g10:universal-dna:gene}{genes} bate com o dos dois
\omterm{def:g10:universal-dna:information}{cromossomos} de chimpanzé \omterm{def:g10:universal-dna:gene}{gene} por \omterm{def:g10:universal-dna:gene}{gene}. Descender de um conjunto comum
de \omterm{def:g10:universal-dna:information}{cromossomos}, modificado uma vez, é o que o parentesco prevê.
\end{solution}

\begin{solution}{exo:g12:human-evolution:7}
O andar ereto, por cerca de dois milhões de anos: pegadas e bacias de
3.6 a 3.2 milhões de anos pertencem a andarilhos com cérebros de
\qty{400}{cm^3}; o crescimento do cérebro começa por volta de 2.5
milhões de anos.
\end{solution}

\begin{solution}{exo:g12:human-evolution:8}
Há dois milhões de anos: \emph{A. africanus} (por pouco),
\emph{Paranthropus}, \emph{Homo habilis} e o \emph{H. erectus} inicial.
Há 100\,000 anos: \emph{H. erectus} (últimas populações),
\emph{H. heidelbergensis} (última), neandertais e \emph{H. sapiens}.
\end{solution}

\begin{solution}{exo:g12:human-evolution:9}
Cruzamentos entre os seres humanos modernos que saíam da África e os
neandertais que eles encontraram na Ásia ocidental e na Europa uns
50\,000 anos atrás. As populações que ficaram na África subsaariana
nunca encontraram neandertais, e assim não carregam nenhum.
\end{solution}

\begin{solution}{exo:g12:human-evolution:10}
Ereto (abertura adiantada), cérebro pequeno, mandíbulas grandes, 2.8
milhões de anos: um australopiteco, colocado num ramo lateral perto da
base da linhagem --- um parente dos ancestrais do \emph{Homo}, e não um
ancestral demonstrado.
\end{solution}

\begin{solution}{exo:g12:human-evolution:11}
Os grandes primatas antropoides formam um \omterm{prop:g12:phylogenetic-trees:rules}{clado}, e os seres humanos são
um ramo dentro dele; já que um \omterm{prop:g12:phylogenetic-trees:rules}{clado} inclui todos os descendentes do nó
dele, os seres humanos são antropoides, como são primatas e mamíferos.
``Descender dos antropoides'' coloca erradamente os antropoides num nó;
os antropoides vivos são pontas irmãs.
\end{solution}

\begin{solution}{exo:g12:human-evolution:12}
Hipótese 1: competição --- os seres humanos modernos, mais numerosos ou
mais bem organizados, tomaram os recursos; a evidência seria um declínio
dos sítios neandertais onde os dois se sobrepuseram. Hipótese 2:
absorção --- populações neandertais pequenas foram em parte absorvidas
pelos cruzamentos e em parte se extinguiram por deriva e por acaso num
clima que mudava; a evidência seria o \omterm{def:g10:universal-dna:information}{DNA} neandertal em nós e sinais de
populações neandertais pequenas e em declínio antes do contato. As duas
podem ser verdadeiras.
\end{solution}

\begin{solution}{exo:g12:human-evolution:13}
Cada etapa da dispersão foi feita por um grupo pequeno que carregava
uma amostra dos \omterm{def:g10:universal-dna:gene}{alelos} da \omterm{def:g12:selection-drift-speciation:population}{população} que ele deixava; cada região nova
foi fundada a partir da anterior. A diversidade caiu a cada fundação, e
assim as populações mais distantes têm a menor. A dispersão se deu por
migrações pequenas e sucessivas, e não por um movimento em massa.
\end{solution}

\begin{solution}{exo:g12:human-evolution:14}
A pecuária leiteira fez do leite um alimento de adultos: os \omterm{def:g11:genes-and-disease:genetic}{portadores}
do \omterm{def:g10:universal-dna:gene}{alelo} da persistência se alimentavam melhor e deixavam mais
descendentes. A migração para latitudes de sol fraco tornou a vitamina
D escassa: a pele mais clara, que a produz em maior quantidade, foi
favorecida. A agricultura criou populações densas e água parada onde a
malária se espalhava: os \omterm{def:g10:universal-dna:gene}{alelos} que protegem contra ela subiram. Em cada
caso, uma mudança no modo de vida mudou quais \omterm{def:g10:universal-dna:gene}{alelos} se reproduziam.
\end{solution}

\begin{solution}{exo:g12:human-evolution:15}
A seleção é qualquer diferença de reprodução entre \omterm{def:g11:genes-and-disease:genetic}{portadores} de \omterm{def:g10:universal-dna:gene}{alelos}
diferentes; a medicina muda quais \omterm{def:g10:universal-dna:gene}{alelos} importam, e não se algum
importa --- a resistência a infecções, a fertilidade, a idade da
reprodução ainda variam e são herdadas. Os três caracteres mostram a
seleção agindo nos últimos dez mil anos, impulsionada pela própria
cultura. E a deriva continua em toda \omterm{def:g12:selection-drift-speciation:population}{população} de qualquer jeito. A
evolução humana mudou de direção, não parou.
\end{solution}

\begin{solution}{pb:g12:human-evolution:1}
\textbf{1.} A: $21.5/0.15 \approx \qty{1.43}{m}$; B: $18.5/0.15
\approx \qty{1.23}{m}$.

\textbf{2.} $0.8 \times 1.43 = \qty{1.14}{m}$ e $0.8 \times 1.23 =
\qty{0.98}{m}$: as duas passadas são mais curtas que a de um andar
lento; eles estavam andando, sem pressa.

\textbf{3.} Um calcanhar fundo: o calcanhar toca primeiro, como no
andar humano; um arco: um pé que age como alavanca e mola; um dedão
paralelo: um pé para impulsionar, e não para agarrar galhos.

\textbf{4.} Um andar de passada larga, calcanhar primeiro, com um pé
parecido com o humano, mantido ao longo de \qty{27}{m}: bipedia
habitual. Desconhecidos a partir das marcas: a bacia, o ângulo do
fêmur, a abertura medular, as curvas da coluna.

\textbf{5.} Que o andar ereto precedeu o aumento do cérebro em mais de
um milhão de anos.

\textbf{6.} 1: australopiteco. 2: \emph{Homo habilis}. 3: \emph{Homo
erectus}. 4: neandertal. 5: \emph{Homo sapiens}.

\textbf{7.} De 450 aos 3.0 milhões a 650 aos 1.9, 950 a 1.0, depois
1450 e 1350: o crescimento mais rápido fica entre 1.9 e 1.0 milhão de
anos (\qty{300}{cm^3} por milhão de anos), depois entre 1.0 e 0.06.

\textbf{8.} Não: os cérebros neandertais eram tão grandes quanto ou
maiores. Os seres humanos modernos são identificados pela testa alta,
pela face pequena sob a caixa craniana, pelo queixo e pelo crânio arredondado.

\textbf{9.} Duas populações do arbusto vivendo lado a lado, distintas
na anatomia, que se cruzaram: o \omterm{def:g10:universal-dna:information}{DNA} neandertal nas pessoas vivas de
fora da África é a evidência.

\textbf{10.} Um fóssil é uma ponta no próprio ramo lateral; a
combinação de caracteres dele o coloca perto do nó de que as \omterm{def:g10:biodiversity-scales:species}{espécies}
posteriores descendem, mas caráter nenhum prova que ele era exatamente
a \omterm{def:g12:selection-drift-speciation:population}{população} que as originou. Ele é um parente próximo do ancestral, que
é tudo o que se consegue demonstrar de um fóssil.

\textbf{11.} Segmentos que perfazem 1.6\% do genoma dela são mais
próximos da sequência neandertal que de qualquer sequência moderna:
eles foram herdados de ancestrais neandertais que se cruzaram com seres
humanos modernos.

\textbf{12.} O gráfico dá cerca de 1.8\% de neandertal na Europa e
cerca de 3.5\% de \omterm{def:g10:universal-dna:information}{DNA} da \omterm{def:g12:selection-drift-speciation:population}{população} asiática na Melanésia: a ascendência
europeia dela trouxe a parcela neandertal, e a ascendência melanésia a
outra (diluída pela descendência mista dela).

\textbf{13.} $50\,000/25 = 2000$ gerações. A parcela não é diluída
porque a \omterm{def:g12:selection-drift-speciation:population}{população} inteira a carrega: misturar-se com outros \omterm{def:g11:genes-and-disease:genetic}{portadores}
mantém a média constante; a diluição só acontece quando os \omterm{def:g11:genes-and-disease:genetic}{portadores}
se misturam com não \omterm{def:g11:genes-and-disease:genetic}{portadores}, e fora da África não havia nenhum.

\textbf{14.} Que a fronteira não estava selada: depois de uns 500\,000
anos de separação, as populações ainda conseguiam produzir descendentes
férteis --- \omterm{def:g10:biodiversity-scales:species}{espécies} em formação, isoladas mais pela geografia que pela
biologia.

\textbf{15.} As \omterm{def:g10:cells-common-unit:organelle}{mitocôndrias} passam só pelas mães: nenhuma linha
materna neandertal sobreviveu entre os seres humanos vivos, seja porque
os descendentes das mães neandertais morreram, seja porque os
cruzamentos se deram sobretudo num sentido só. Os traços nucleares
sobreviveram; aquela linhagem não.

\textbf{16.} $1.2/0.27 \approx 4.4$ milhões de anos --- uma estimativa
por baixo, já que o relógio foi ajustado em nós mais distantes.

\textbf{17.} Os fósseis são mais antigos que a estimativa do relógio; o
nó verdadeiro tem pelo menos a idade do fóssil mais antigo da linhagem,
ou seja, de 6 a 7 milhões de anos, e a taxa do relógio para esse
intervalo recente deve ser um pouco mais lenta que 0.27\%.

\textbf{18.} O \emph{H. erectus} e, na linha do tempo, os últimos
\emph{Paranthropus} (e, a partir de cerca de 0.7, o
\emph{H. heidelbergensis}); desses, só a linhagem que passa pelo
\emph{H. heidelbergensis} até o \emph{H. sapiens} tem descendentes
vivos, e o \emph{H. erectus} apenas através de populações que
desaguaram nela --- uma linha entre várias.

\textbf{19.} O esqueleto e o cérebro mudaram pouco nesse tempo,
enquanto as ferramentas, a arte, a linguagem, a agricultura e as
cidades transformaram o modo como os seres humanos vivem; as mudanças
são transmitidas pelo aprendizado, se acumulam dentro das gerações e
hoje moldam o ambiente em que a \omterm{def:g10:biodiversity-scales:species}{espécie} evolui.

\textbf{20.} O caminhante A media cerca de \qty{1.4}{m} e andava
devagar sobre pés parecidos com os humanos; os pés vieram primeiro, o
cérebro um milhão de anos depois, o queixo por último; a estudante
carrega \omterm{def:g10:universal-dna:gene}{genes} de duas populações extintas.
\end{solution}
